\documentclass[12pt, a4paper]{report}
\usepackage{fontspec}
\setmainfont{CMU Serif}
\usepackage{marginnote}
\usepackage[german]{babel}
\usepackage{xurl}
\usepackage[colorlinks=true, linkcolor=blue, urlcolor=blue]{hyperref}
%Section auf größe von subsection setzen
\usepackage{titlesec}
\titleformat*{\section}{\large\bfseries}
%Grafik
\usepackage{float}
%\usepackage{pgf-pie}
\usepackage{pdfpages}
%\usepackage{pgf, tikz} %Tikz GRAPHICS

%Seitenlayout, Margins verkleinern
\usepackage[a4paper, headheight=1.75cm, top=3.2cm,bottom=30mm, left=4.5cm, right=1.5cm, marginparwidth=4cm, marginparsep=.4cm]{geometry} 
\usepackage{setspace}%Zeilenabstand verändern zB \onehalfspacing(Präambel)

\usepackage{amsmath}
\usepackage{csquotes}
\usepackage{graphicx} %Für Bilder
\usepackage{subcaption} %Für subfigure
\usepackage{siunitx}
%Bruchstrich
\sisetup{per-mode=fraction, output-decimal-marker = {,}, exponent-product=\cdot}
%Hier beliebige Einheiten hinzufügen
\DeclareSIUnit\year{Jahre}

%Kopf-Fußzeile
\usepackage{fancyhdr}
%\fancyhf{}
\pagestyle{fancy}
\fancyhead[L]{\Large{\textsc{Jahresprogramm}}}
\fancyhead[R]{\includegraphics[width=1.75cm]{logo.png}}
\renewcommand{\headrulewidth}{0.5pt}
\setlength{\headheight}{56.4pt}
\cfoot{\thepage}

%Paragraph Settings
\setlength{\parindent}{0cm}

%Titel
\title{Vorläufiges Fahrtenprogramm}
\date{}
%\author{Leonhard Schraut}
%Dokument

%%%%%%%%%%%%%%%%%%%%% Generated File %%%%%%%%%%%%%%%%%%%%%
\fancypagestyle{Allgemein}{
    \fancyhead[L]{\Large{\textsc{Allgemein}}}
    \fancyhead[R]{\includegraphics[width=1.75cm]{logo.png}}
    \renewcommand{\headrulewidth}{0.5pt}
    \cfoot{\thepage}
}

\fancypagestyle{Jugend}{
    \fancyhead[L]{\Large{\textsc{Jugend}}}
    \fancyhead[R]{\includegraphics[width=1.75cm]{logo.png}}
    \renewcommand{\headrulewidth}{0.5pt}
    \cfoot{\thepage}
}

\fancypagestyle{Kanuwandern}{
    \fancyhead[L]{\Large{\textsc{Kanuwandern}}}
    \fancyhead[R]{\includegraphics[width=1.75cm]{logo.png}}
    \renewcommand{\headrulewidth}{0.5pt}
    \cfoot{\thepage}
}

\fancypagestyle{Wildwasser}{
    \fancyhead[L]{\Large{\textsc{Wildwasser}}}
    \fancyhead[R]{\includegraphics[width=1.75cm]{logo.png}}
    \renewcommand{\headrulewidth}{0.5pt}
    \cfoot{\thepage}
}

\fancypagestyle{Kanupolo}{
    \fancyhead[L]{\Large{\textsc{Kanupolo}}}
    \fancyhead[R]{\includegraphics[width=1.75cm]{logo.png}}
    \renewcommand{\headrulewidth}{0.5pt}
    \cfoot{\thepage}
}

\begin{document}

\begingroup
    \hypersetup{hidelinks}
    \tableofcontents\thispagestyle{fancy}
\endgroup
\reversemarginpar

\section*{Bemerkungen}\paragraph{Allgemein}\mbox{}\\
Alle Termine sind bald auch in Klubraum zu finden.
\paragraph{Wildwasser}\mbox{}\\
Die Teilnahme an allen Fahrten erfolgt auf eigene Gefahr. Jede Haftung des Veranstalters und seiner Hilfspersonen für Personen- und Sachschäden ist ausgeschlossen, außer bei Vorsatz oder grober Fahrlässigkeit. Auch bei Wanderfahrten ist das Tragen von Schwimmwesten Pflicht. Ebenso ist auf, der Wassertemperatur angepasster, Kleidung zu achten.


Wildwasserfahrten nur mit entsprechender Ausrüstung: Wildwasserboot mit Auftriebskörpern und Bergegriffen, Schutzhelm, Wildwasserschwimmweste, Wurfsack, Kälteschutzkleidung für das Wasser (Neopren, Paddeljacke oder Trockenanzug), feste Paddelschuhe. Eine Haftung für Sachschäden oder Diebstahl ist ausgeschlossen.

\paragraph{Kanuwandern}\mbox{}\\
Zu beziehungsweise zwischen den Terminen werden natürlich noch weitere Fahrten stattfinden untder dem Motto, ,,wer hat Lust'' oder ,,wir könnten doch mal gemeinsam auf dem Fluss oder dem Bach fahren''.

Für alle Termine gibt es Infos am schwarzen Brett oder auf der Homepage des KC-Würzburgs.

Informationen zu Schulungen und Veranstaltungen: \url{www.kanu-unterfranken.de}, bzw \url{www.kanu-bayern.de}, \url{www.kanu.de}.

Bitte führt ein Fahrtenbuch, das ist per EFB oder auf papier möglich. Das Fahrtenbuch ist für den persönlichen Versicherungsschutz und für den Verein bzw. den Verband von Vorteil.\chapter*{Allgemein}
\thispagestyle{Allgemein}
\addcontentsline{toc}{chapter}{\protect\numberline{}Allgemein}
\pagestyle{Allgemein}
\section*{März}\paragraph{Arbeitsdienst}\marginnote{25.03 \\ ab 09:00}
\begin{itemize}
    \item Ansprechpartner: Vorstand \href{mailto:vorstand@kc-wuerzburg.de}{vorstand@kc-wuerzburg.de}
    \item Arbeitsdienst am KCW
\end{itemize}

\section*{Mai}\paragraph{Besuch im Walderlebniszentrum Einsiedel}\marginnote{13.05}
\begin{itemize}
    \item Ansprechpartner: Michael \href{mailto:vergnuegung@kc-wuerzburg.de}{vergnuegung@kc-wuerzburg.de}
    \item Startzeit und Treffpunkt werden noch bekannt gegeben
\end{itemize}

\section*{Juli}\paragraph{Sommerfest}\marginnote{01.07 \\ ab 17:00}
\begin{itemize}
    \item Ansprechpartner: Michael \href{mailto:vergnuegung@kc-wuerzburg.de}{vergnuegung@kc-wuerzburg.de}
    \item Sommerfest am Bootssteg
\end{itemize}

\section*{Oktober}\paragraph{Besuch im Röntgen-Museum}\marginnote{07.10}
\begin{itemize}
    \item Ansprechpartner: Michael \href{mailto:vergnuegung@kc-wuerzburg.de}{vergnuegung@kc-wuerzburg.de}
    \item Startzeit und Treffpunkt werden noch bekannt gegeben
\end{itemize}

\paragraph{Jahreshauptversammlung}\marginnote{13.10 \\ ab 19:30}
\begin{itemize}
    \item Ansprechpartner: Vorstand \href{mailto:vorstand@kc-wuerzburg.de}{vorstand@kc-wuerzburg.de}
    \item Ort: Gaststätte KCW
    \item Mit Neuwahlen
\end{itemize}

\section*{November}\paragraph{Arbeitsdienst}\marginnote{18.11 \\ ab 09:00}
\begin{itemize}
    \item Ansprechpartner: Vorstand \href{mailto:vorstand@kc-wuerzburg.de}{vorstand@kc-wuerzburg.de}
    \item Arbeitsdienst am KCW
\end{itemize}

\section*{Dezember}\paragraph{Nikolaus am Bootssteg}\marginnote{10.12 \\ ab 15:00}
\begin{itemize}
    \item Ansprechpartner: Michael \href{mailto:vergnuegung@kc-wuerzburg.de}{vergnuegung@kc-wuerzburg.de}
\end{itemize}

\paragraph{Weihnachtsfeier mit Programm}\marginnote{16.12 \\ ab 19:00}
\begin{itemize}
    \item Ansprechpartner: Michael \href{mailto:vergnuegung@kc-wuerzburg.de}{vergnuegung@kc-wuerzburg.de}
    \item Ort: Vereinsgaststätte
\end{itemize}


\chapter*{Jugend}
\thispagestyle{Jugend}
\addcontentsline{toc}{chapter}{\protect\numberline{}Jugend}
\pagestyle{Jugend}
\section*{April}\paragraph{Huningue}\marginnote{15.04 - 16.04 \\ ab 06:00}
\begin{itemize}
    \item Ansprechpartnerin: Julia \href{mailto:jugend@kc-wuerzburg.de}{jugend@kc-wuerzburg.de}
    \item Ort: Huningue
    \item Übernachtung im gemietetem Haus
    \item Wildwassertraining auf dem Kanal
    \item Vor allem für Jugend und Anfänger gedacht
\end{itemize}

\section*{August}\paragraph{Fun-Turnier Jena}\marginnote{26.08 - 27.08}
\begin{itemize}
    \item Ansprechpartner: Leo \href{mailto:sport@kc-wuerzburg.de}{sport@kc-wuerzburg.de}
    \item Auch für Turnierneulinge geeignet
\end{itemize}

\section*{September}\paragraph{BKV Lehrgang Lechtal}\marginnote{16.09 - 17.09}
\begin{itemize}
    \item Ansprechpartner: Bernd Sachs \href{mailto:ww-breitensport@kanu-bayern.de}{ww-breitensport@kanu-bayern.de}
    \item Ansprechpartner KCW: Leo \href{mailto:sport@kc-wuerzburg.de}{sport@kc-wuerzburg.de}
    \item Schulung für Erwachsene sowie Familien (incl. Kinder)
    \item Schwierigkeit: WW I bis WW III
    \item Ort: Vorderhornbach
    \item Weitere Informationen und Anmeldung: \url{www.kanu-bayern.de/Freizeitsport/WW-Breitensport/Termine/644/Wildwasser-Lehrgang-3-Schulung-fuer-Erwachsene-sowie-Familien-incl-Kinder/}
\end{itemize}


\chapter*{Kanuwandern}
\thispagestyle{Kanuwandern}
\addcontentsline{toc}{chapter}{\protect\numberline{}Kanuwandern}
\pagestyle{Kanuwandern}
\section*{April}\paragraph{Karfreitags Paddeln}\marginnote{07.04 \\ ab 08:30}
\begin{itemize}
    \item Ansprechpartner: Rainer \href{mailto:kanuwandern@kc-wuerzburg.de}{kanuwandern@kc-wuerzburg.de}
    \item Traditionelles, leicht sportliches ,,Karfreitags-Paddeln'' auf einem Kleinfluss in der Nähe
    \item Treffpunkt: am KCW
\end{itemize}

\paragraph{Anpaddeln}\marginnote{23.04 \\ ab 10:00}
\begin{itemize}
    \item Ansprechpartner: Rainer \href{mailto:kanuwandern@kc-wuerzburg.de}{kanuwandern@kc-wuerzburg.de}
    \item Strecke: Ochsenfurt bis zum KCW (ca. 20km)
    \item Kurzstrecke: Winterhausen zum KCW (ca. 12km)
    \item Treffpunkt: am KCW
\end{itemize}

\paragraph{Bezirkswanderfahrt in Würzburg}\marginnote{28.04 - 30.04 \\ ab 15:00}
\begin{itemize}
    \item Ansprechpartner: Edgar König \href{mailto:kanuwandern@kanu-unterfranken.de}{kanuwandern@kanu-unterfranken.de}
    \item Ansprechpartner KCW: Rainer \href{mailto:kanuwandern@kc-wuerzburg.de}{kanuwandern@kc-wuerzburg.de}
    \item Wanderfahrten von ca. 20km Länge
    \item Anmeldung ung Gebühren ca. 2 Wochen vorher
\end{itemize}

\section*{Juni}\paragraph{50. Bayrisches Wanderfahrertreffen im ,,Herrgottswinkel''}\marginnote{14.06 - 18.06}
\begin{itemize}
    \item Ansprechpartner: Rainer \href{mailto:kanuwandern@kc-wuerzburg.de}{kanuwandern@kc-wuerzburg.de}
    \item Anmeldung: \url{www.kanu-bayern.de/Freizeitsport/Kanuwandern/Termine/668/50-Bayerisches-Kanu-Wanderfahrertreffen/}
    \item Veranstalter: TSV Schongau
\end{itemize}

\section*{Juli}\paragraph{Frauenfahrt}\marginnote{07.07 - 09.07}
\begin{itemize}
    \item Ansprechpartnerin: Angelika Gardemann \href{mailto:frauen@kanu-unterfranken.de}{frauen@kanu-unterfranken.de}
    \item Ansprechpartner KCW: Rainer \href{mailto:kanuwandern@kc-wuerzburg.de}{kanuwandern@kc-wuerzburg.de}
    \item Kanufrauen Touren auf dem Main
    \item Ausrichter: TG Würzburg-Heidingsfeld
\end{itemize}

\section*{August}\paragraph{Deutsch, italienisches, niederländisches Freundschaftstreffen}\marginnote{07.08 - 10.08}
\begin{itemize}
    \item Ansprechpartner: Rainer \href{mailto:kanuwandern@kc-wuerzburg.de}{kanuwandern@kc-wuerzburg.de}
    \item Ausrichter: KC-Klingenberg
\end{itemize}

\paragraph{Urlaubsfahrt der Unterfranken}\marginnote{27.08 - 02.09}
\begin{itemize}
    \item Ansprechpartner: Edgar König \href{mailto:kanuwandern@kanu-unterfranken.de}{kanuwandern@kanu-unterfranken.de}
    \item Ansprechpartner KCW: Rainer \href{mailto:kanuwandern@kc-wuerzburg.de}{kanuwandern@kc-wuerzburg.de}
    \item Ziel: Labusee
\end{itemize}

\section*{September}\paragraph{Zwetschgenkuchenfahrt}\marginnote{23.09 \\ ab 08:00}
\begin{itemize}
    \item Ansprechpartner: Rainer \href{mailto:kanuwandern@kc-wuerzburg.de}{kanuwandern@kc-wuerzburg.de}
    \item Strecke: Mainschleife von Volkach bis Albertshofen
    \item Meldeschluss: 11.09
    \item Anmeldunng via Schwarzes Brett
    \item Bootsverladung: Freitag, ca 17:00
\end{itemize}

\paragraph{Abpaddeln}\marginnote{24.09 \\ ab 10:00}
\begin{itemize}
    \item Ansprechpartner: Rainer \href{mailto:kanuwandern@kc-wuerzburg.de}{kanuwandern@kc-wuerzburg.de}
    \item Strecke: Ochsenfurt bis KC Würzburg (ca. 20km)
    \item Kurzstrecke: Winterhausen bis KC Würzbg (ca. 12km)
    \item Treffpunkt: am KCW
\end{itemize}


\chapter*{Wildwasser}
\thispagestyle{Wildwasser}
\addcontentsline{toc}{chapter}{\protect\numberline{}Wildwasser}
\pagestyle{Wildwasser}
\section*{März}\paragraph{44. Wildwasserfahrt auf erzgebirgischen Wildflüssen}\marginnote{31.03 - 02.04}
\begin{itemize}
    \item Ansprechpartnerin: Anne \href{mailto:anne\_fathke@arcor.de}{anne\_fathke@arcor.de}
    \item Ansprechpartner KCW: Volker \href{mailto:v.woehrn@arcor.de}{v.woehrn@arcor.de}
    \item Ort:  Warmbad Wolkenstein
    \item Ausrichter: Sportclub Berlin - Grünau e.V.
    \item Anmeldung bereits abgeschlossen, Nachmeldung evtl. auf Nachfrage möglich
\end{itemize}

\section*{April}\paragraph{Huningue}\marginnote{15.04 - 16.04 \\ ab 06:00}
\begin{itemize}
    \item Ansprechpartner: Sebastian \href{mailto:wildwasser@kc-wuerzburg.de}{wildwasser@kc-wuerzburg.de}
    \item Ort: Huningue
    \item Übernachtung im gemietetem Haus
    \item Wildwassertraining auf dem Kanal
    \item Vor allem für Jugend und Anfänger gedacht
\end{itemize}

\paragraph{Erster Mai Fahrt}\marginnote{28.04 - 01.05}
\begin{itemize}
    \item Ansprechpartner: Sebastian \href{mailto:wildwasser@kc-wuerzburg.de}{wildwasser@kc-wuerzburg.de}
    \item Schwierigkeit: WW II bis WW III
    \item Ort: Häselgehr
    \item Gepaddelt werden: Lech, Loisach, Ammer
\end{itemize}

\section*{Mai}\paragraph{BKV Lehrgang Wildalpen}\marginnote{18.05 - 21.05}
\begin{itemize}
    \item Ansprechpartner: Bernd Sachs \href{mailto:ww-breitensport@kanu-bayern.de}{ww-breitensport@kanu-bayern.de}
    \item Ansprechpartner KCW: Sebastian \href{mailto:wildwasser@kc-wuerzburg.de}{wildwasser@kc-wuerzburg.de}
    \item Schwierigkeit: WW I bis WW III
    \item Ort: Wildalpen
    \item Meldeschluss: 21.04.2023 (sehr schnell ausgebucht)
    \item Anmeldung: \url{www.kanu-bayern.de/Freizeitsport/WW-Breitensport/Termine/642/Wildwasser-Lehrgang-1/}
\end{itemize}

\paragraph{Pfingst Fahrt}\marginnote{25.05 - 29.05}
\begin{itemize}
    \item Ansprechpartner: Sebastian \href{mailto:wildwasser@kc-wuerzburg.de}{wildwasser@kc-wuerzburg.de}
    \item Ort: noch unbekannt
\end{itemize}

\section*{Juni}\paragraph{Durance/Seealpen}\marginnote{08.06 - 18.06}
\begin{itemize}
    \item Ansprechpartner: Sebastian \href{mailto:wildwasser@kc-wuerzburg.de}{wildwasser@kc-wuerzburg.de}
    \item Schwierigkeit: WW II bis WW IV
    \item Bis 13.06 Durance, anschließend weiterfahrt z den Seealpen
\end{itemize}

\section*{Juli}\paragraph{BKV Lehrgang Augsburg}\marginnote{14.07 - 16.07}
\begin{itemize}
    \item Ansprechpartner: Bernd Sachs \href{mailto:ww-breitensport@kanu-bayern.de}{ww-breitensport@kanu-bayern.de}
    \item Ansprechpartner KCW: Sebastian \href{mailto:wildwasser@kc-wuerzburg.de}{wildwasser@kc-wuerzburg.de}
    \item Technik-Lehrgang am Augsburger Eiskanal
    \item Weitere Informationen und Anmeldung: \url{www.kanu-bayern.de/Freizeitsport/WW-Breitensport/Termine/643/Wildwasser-Lehrgang-2/}
\end{itemize}

\paragraph{Plattling}\marginnote{29.07 - 30.07}
\begin{itemize}
    \item Ansprechpartner: Leo \href{mailto:sport@kc-wuerzburg.de}{sport@kc-wuerzburg.de}
    \item Surfen in der Welle in Plattling
    \item Abhängig vom Wasserstand
\end{itemize}

\section*{September}\paragraph{BKV Lehrgang Lechtal}\marginnote{16.09 - 17.09}
\begin{itemize}
    \item Ansprechpartner: Bernd Sachs \href{mailto:ww-breitensport@kanu-bayern.de}{ww-breitensport@kanu-bayern.de}
    \item Ansprechpartner KCW: Sebastian \href{mailto:wildwasser@kc-wuerzburg.de}{wildwasser@kc-wuerzburg.de}
    \item Schulung für Erwachsene sowie Familien (incl. Kinder)
    \item Schwierigkeit: WW I bis WW III
    \item Ort: Vorderhornbach
    \item Weitere Informationen und Anmeldung: \url{www.kanu-bayern.de/Freizeitsport/WW-Breitensport/Termine/644/Wildwasser-Lehrgang-3-Schulung-fuer-Erwachsene-sowie-Familien-incl-Kinder/}
\end{itemize}

\paragraph{Haiming}\marginnote{29.09 - 03.10}
\begin{itemize}
    \item Ansprechpartner: Sebastian \href{mailto:wildwasser@kc-wuerzburg.de}{wildwasser@kc-wuerzburg.de}
    \item Schwierigkeit: WW II bis WW IV
    \item Ort: Haiming
    \item Gepaddelt werden: Ötz, Inn und Sanna
\end{itemize}


\chapter*{Kanupolo}
\thispagestyle{Kanupolo}
\addcontentsline{toc}{chapter}{\protect\numberline{}Kanupolo}
\pagestyle{Kanupolo}
\section*{Mai}\paragraph{Trainingstag Kanupolo}\marginnote{06.05 \\ ab 09:00}
\begin{itemize}
    \item Ansprechpartner: Timo \href{mailto:kanupolo@kc-wuerzburg.de}{kanupolo@kc-wuerzburg.de}
    \item Trainingstag am KCW
\end{itemize}

\section*{Juni}\paragraph{Trainingstag Kanupolo}\marginnote{03.06 \\ ab 09:00}
\begin{itemize}
    \item Ansprechpartner: Timo \href{mailto:kanupolo@kc-wuerzburg.de}{kanupolo@kc-wuerzburg.de}
    \item Trainingstag am KCW
\end{itemize}

\paragraph{Bayrische Meisterschaft}\marginnote{24.06 - 25.06}
\begin{itemize}
    \item Ansprechpartner: Timo \href{mailto:kanupolo@kc-wuerzburg.de}{kanupolo@kc-wuerzburg.de}
    \item Ausrichter: Kehlheim
\end{itemize}

\section*{August}\paragraph{Fun-Turnier Jena}\marginnote{26.08 - 27.08}
\begin{itemize}
    \item Ansprechpartner: Timo \href{mailto:kanupolo@kc-wuerzburg.de}{kanupolo@kc-wuerzburg.de}
    \item Auch für Turnierneulinge geeignet
\end{itemize}

\section*{September}\paragraph{Kanupolo Turnier Würzburg}\marginnote{01.09 - 03.09}
\begin{itemize}
    \item Ansprechpartner: Leo \href{mailto:sport@kc-wuerzburg.de}{sport@kc-wuerzburg.de}
    \item Termin unter Vorbehalt
\end{itemize}


\end{document}